% BHU PYQ -- Manually transcribed & error-corrected
\documentclass[11pt,a4paper]{article}

\usepackage[a4paper,top=2.5cm,bottom=2.5cm,left=2.8cm,right=2.8cm,headheight=14pt]{geometry}
\usepackage{amsmath,amssymb}
\usepackage[T1]{fontenc}
\usepackage[utf8]{inputenc}
\usepackage{lmodern}
\usepackage{microtype}
\usepackage{fancyhdr}
\usepackage{enumitem}
\usepackage{hyperref}
\usepackage{lastpage}
\usepackage{parskip}

\hypersetup{colorlinks=true,urlcolor=black,linkcolor=black,
  pdftitle={BPT-101: Mechanics and Relativity -- BHU PYQ},pdfauthor={Banaras Hindu University}}

\pagestyle{fancy}\fancyhf{}
\renewcommand{\headrulewidth}{0.4pt}\renewcommand{\footrulewidth}{0.4pt}
\fancyhead[L]{\small   \textit{Banaras Hindu University}}
\fancyhead[R]{\small   \textit{BPT-101: Mechanics and Relativity}}
\fancyfoot[C]{\small Page  \thepage\ of \pageref{LastPage}}


\newlist{parts}{enumerate}{2}
\setlist[parts,1]{label=(\alph*),leftmargin=*,itemsep=5pt,topsep=3pt}
\setlist[parts,2]{label=(\roman*),leftmargin=*,itemsep=3pt,topsep=2pt}

\newcommand{\pts}[1]{\hfill   \textit{[#1]}}


\begin{document}


\begin{center}
  {\large\bfseries BANARAS HINDU UNIVERSITY}\\[2pt]
  {
\normalsize B.Sc. (Hons.) Semester I Examination, 2016-17}\\[6pt]
  \rule{0.8\linewidth}{0.5pt}\\[6pt]
  {\large\bfseries Paper No. BPT-101: Mechanics and Relativity}\\[4pt]
  {
\normalsize Time: 3 Hours\hfill Maximum Marks: 70}
\end{center}
\rule{\linewidth}{0.4pt}
\smallskip



\noindent   \textit{Instructions: Answer all questions. Symbols have their usual meanings.}
\bigskip




\noindent   \textbf{Question 1.} Answer any   \textbf{four} of the following: \pts{4 $ \times$ 3.5 = 14}

\begin{parts}
  \item What are fictitious (pseudo) forces and when do they come into play? Calculate the fictitious force and observed (total) force on a body of mass $5\, \text{kg}$ in a frame of reference moving:
    
\begin{parts}
      \item vertically upwards with an acceleration of $4\, \text{m/s}^2$.
      \item vertically downwards with an acceleration of $4\, \text{m/s}^2$.
    \end{parts}
    (Take $g = 9.8\, \text{m/s}^2$).
  \item What is the rest mass of a photon? Give arguments in support of your answer.
  \item What are central forces? Show that the angular momentum of a particle moving under the influence of a central force always remains constant.
  \item In a nuclear reactor, hydrogen (in the form of water or heavy water) is used as a moderator. Give reasons.
  \item Obtain an expression for the time period of a compound pendulum and discuss its advantages over a simple pendulum.
  \item What are damped and undamped oscillations? Discuss with examples. In an oscillatory LCR circuit, $L = 0.2\, \text{H}$ and $C = 1.2\,\mu \text{F}$. What is the maximum value of resistance $R$ for the circuit to remain oscillatory?
\end{parts}

\medskip\rule{\linewidth}{0.2pt}




\noindent   \textbf{Question 2.}

\begin{parts}
  \item What is a cantilever? A cantilever of rectangular cross-section is loaded at its free end. Obtain an expression for the depression of the cantilever at its free end. \pts{9}
  \item Calculate Poisson's ratio for silver. Given that Young's modulus for silver is $7.25  \times 10^{10}\, \text{N/m}^2$ and its bulk modulus is $1.1  \times 10^{11}\, \text{N/m}^2$. \pts{5}
\end{parts}

\medskip\rule{\linewidth}{0.2pt}




\noindent   \textbf{Question 3.}

\begin{parts}
  \item On the basis of the Lorentz transformation equations, discuss:
    
\begin{parts}
      \item length contraction.
      \item time dilation.
    \end{parts} \pts{10}
  \item A proton of rest mass $1.67  \times 10^{-27}\, \text{kg}$ is moving with a velocity of $0.9c$. Find its relativistic mass and momentum. \pts{4}
\end{parts}
\hfill   \textbf{OR}

\begin{parts}
  \item What are the assumptions in Poiseuille's equation for the flow of a liquid through capillary tubes? Show that the coefficient of viscosity is given by:
    \[ \eta = \frac{\pi P r^4}{8 V L} \]
    where $V$ is the volume of the liquid flowing per unit time and the other symbols have their usual meanings. \pts{8}
  \item Explain how you will determine the coefficient of viscosity of water by Poiseuille's method in the laboratory. Discuss the effect of temperature on the viscosity of liquids. \pts{6}
\end{parts}

\medskip\rule{\linewidth}{0.2pt}




\noindent   \textbf{Question 4.}

\begin{parts}
  \item What do you mean by relativistic mass? Derive the relativistic mass-energy relation $E = mc^2$ and discuss its physical importance. \pts{8}
  \item If the relativistic kinetic energy is given by $K = mc^2 - m_0 c^2$, show that its classical value (when $v \ll c$) reduces to $\frac{1}{2} m_0 v^2$. \pts{6}
\end{parts}
\hfill   \textbf{OR}

\begin{parts}
  \item Derive an expression for the twisting couple per unit angular twist for a solid cylinder. Show that for cylinders of the same mass, material, and length, the twisting couple per unit twist is greater for a hollow cylinder than for a solid one. \pts{8}
  \item Obtain an expression for the horizontal deflection of a body falling freely under the action of gravity at a latitude $\lambda$ due to the Coriolis force. \pts{6}
\end{parts}

\medskip\rule{\linewidth}{0.2pt}




\noindent   \textbf{Question 5.}

\begin{parts}
  \item Calculate the displacement suffered by a stone dropped from the top of a tower $100\, \text{m}$ high due to the Coriolis force:
    
\begin{parts}
      \item in latitude $60^\circ \text{N}$.
      \item at the equator.
    \end{parts}
    (Take angular velocity of the Earth $\omega = 7.29  \times 10^{-5}\, \text{rad/s}$ and $g = 9.8\, \text{m/s}^2$). \pts{14}
\end{parts}
\hfill   \textbf{OR}

\begin{parts}
  \item What do you understand by the laboratory (L) and centre of mass (C) frames of reference? Discuss two-dimensional elastic collisions in both frames of reference. Discuss the advantages of studying collision processes in the centre of mass frame of reference. \pts{14}
\end{parts}

\vfill
\rule{\linewidth}{0.4pt}

\begin{center}\small
  \href{https://bhu-exam.vercel.app/}{Click here to visit our website}\\[4pt]
  \href{https://me-aryan.vercel.app/}{Click here to visit our contributor}\\[4pt]
  \href{https://quashan-me.vercel.app/}{Click here to visit our contributor}
\end{center}

\end{document}
